\documentclass[11pt,a4paper]{article}
\usepackage[utf8]{inputenc}
\usepackage[spanish]{babel}
\usepackage{amsmath}
\usepackage{amsfonts}
\usepackage{amssymb}
\usepackage{geometry}
\usepackage{booktabs}
\usepackage{hyperref}

\geometry{left=2.5cm, right=2.5cm, top=2.5cm, bottom=2.5cm}

\title{\textbf{Desarrollo del Modelo Matemático de la Duna Bidispersa\\[4pt]
con Deposición por Avalanchas y Capa Activa Confinada\\[6pt]
\large (\texttt{dune\_bidispersa\_deposicion\_ext.py})}}
\author{Felipe Espinoza}
\date{\today}

\newcommand{\dd}{\mathrm{d}}
\newcommand{\pd}[2]{\dfrac{\partial #1}{\partial #2}}

\begin{document}

\maketitle
\tableofcontents
\newpage

% ======================================================================
\section{Introducción y alcance}
Este documento describe la formulación matemática y numérica del modelo
implementado en \texttt{dune\_bidispersa\_deposicion\_ext.py}. El código
simula la evolución de la concentración volumétrica de finos
$\phi_s(x,\eta,t)\in[0,1]$ dentro de una duna bidispersa (dos tamaños de
grano: finos $d_s$ y gruesos $d_l>d_s$) que migra sobre un lecho. El modelo
parte del esquema clásico de advección--segregación--difusión de Trewhela,
Ancey \& Gray (2021) y de la fenomenología de \emph{grain-flows} de
Kleinhans (2004), pero incorpora cinco adaptaciones que son el objeto
principal de este análisis:
\begin{enumerate}
    \item \textbf{Cambio de variable} $Q=h\,\phi_s$ y transformación a
    coordenadas $\sigma$ (Sección~\ref{sec:cambio}).
    \item \textbf{Confinamiento de la capa activa} y \textbf{desacople} de
    la tasa de corte respecto de la advección (Sección~\ref{sec:confin}).
    \item \textbf{Piso de presión} que regulariza la velocidad de
    segregación (Sección~\ref{sec:trewhela}).
    \item \textbf{Deposición por avalanchas discretas} en la superficie del
    sotavento (Sección~\ref{sec:borde}).
    \item \textbf{Esquema de alto orden} WENO5 (horizontal) y MUSCL+Rusanov
    (vertical) para preservar la laminación (Sección~\ref{sec:numerico}).
\end{enumerate}
La relación de tamaños es $R = d_l/d_s = 3.33$ con $d_l=1.0$~mm y
$d_s=0.3$~mm, y la composición media es constante, $\phi_s=0.70$.

% ======================================================================
\section{Marco comóvil y geometría de la duna}
\label{sec:geom}
La duna migra a celeridad $c_{\mathrm{mig}}$ en el marco de laboratorio.
Para trabajar sobre una malla fija se adopta el \textbf{marco comóvil},
donde la forma de la duna es estacionaria:
\begin{equation}
    x_{\mathrm{lab}} = x + c_{\mathrm{mig}}\,t .
\end{equation}
La topografía triangular (versión \texttt{ext}) se define a trozos:
\begin{equation}
    h(x) = \begin{cases}
        H_{\mathrm{base}} + H_d\,\dfrac{x}{x_{\mathrm{crest}}}, & 0\le x \le x_{\mathrm{crest}} \quad\text{(barlovento / stoss)}\\[10pt]
        H_{\mathrm{base}} + H_d\left(1 - \dfrac{x-x_{\mathrm{crest}}}{L_{\mathrm{lee}}}\right), & x_{\mathrm{crest}} < x \le L_{\mathrm{dune}} \quad\text{(sotavento / lee)},
    \end{cases}
    \label{eq:h_tri}
\end{equation}
con pendiente constante por tramos
\begin{equation}
    \frac{\dd h}{\dd x} = \begin{cases}
        +H_d/x_{\mathrm{crest}} > 0 & \text{(stoss, suave)}\\[4pt]
        -H_d/L_{\mathrm{lee}} < 0 & \text{(lee, al ángulo de reposo }\alpha_{\mathrm{lee}}=30^\circ).
    \end{cases}
\end{equation}
Aquí $H_d=8$~mm es la altura, $L_{\mathrm{dune}}=100$~mm la longitud,
$H_{\mathrm{base}}=1$~mm un lecho base (escala de $\eta$), y
$L_{\mathrm{lee}}=H_d/\tan\alpha_{\mathrm{lee}}$. En la
Sección~\ref{sec:curva} se reemplaza esta geometría por una versión curva
con stoss convexo.

\paragraph{Escala de migración de demostración.}
La celeridad física es $c_{\mathrm{mig,fis}}=0.3$~mm/min. Como a esa
velocidad la duna casi no renueva su volumen en tiempo de cómputo
razonable, se usa un factor de demostración
$c_{\mathrm{mig}} = \texttt{demo\_speedup}\cdot c_{\mathrm{mig,fis}}$ con
$\texttt{demo\_speedup}=10$. Esto \emph{no} altera la física del sorting;
solo acelera el reciclaje del material.

% ======================================================================
\section{Cambio de variable $Q=h\,\phi_s$ y ecuación gobernante}
\label{sec:cambio}
\subsection{Coordenadas $\sigma$ y masa conservada}
Se introduce la coordenada vertical normalizada
\begin{equation}
    \eta = \frac{z}{h(x)} \in [0,1], \qquad z=\eta\,h(x),
\end{equation}
que mapea el dominio físico de altura variable a un rectángulo fijo
($\eta=0$ base, $\eta=1$ superficie). La masa total de finos es
\begin{equation}
    M_{\mathrm{finos}} = \int_0^L\!\!\int_0^{h(x)} \phi_s\,\dd z\,\dd x
    = \int_0^L\!\!\int_0^1 \underbrace{h(x)\,\phi_s}_{\textstyle Q}\,\dd\eta\,\dd x .
\end{equation}
Al resolver directamente para $Q=h\,\phi_s$ (espesor equivalente de finos),
la masa de cada celda es $Q\,\dd x\,\dd\eta$ con $\dd x,\dd\eta$
constantes, lo que garantiza conservación estricta de masa ante cualquier
deformación de $h$.

\subsection{Forma conservativa en $\sigma$}
Partiendo de la conservación cartesiana de la fase fina para un flujo
incompresible ($\nabla\cdot\vec u=0$),
\begin{equation}
    \pd{\phi_s}{t} + \pd{(u\phi_s)}{x} + \pd{(w\phi_s)}{z}
    + \pd{F_{\mathrm{seg}}}{z}
    = \pd{}{z}\!\left(D_{sl}\,\pd{\phi_s}{z}\right),
    \label{eq:cart}
\end{equation}
y aplicando la regla de la cadena a $(x,z,t)\to(x,\eta,t)$ junto con el
cambio $Q=h\phi_s$ (derivación completa en la Sección~\ref{sec:derivacion}),
se obtiene la \textbf{forma conservativa estricta} que resuelve el código:
\begin{equation}
    \boxed{\;\pd{Q}{t} + \pd{(u\,Q)}{x}
    + \pd{}{\eta}\Big[\, F_{\mathrm{adv},\eta} + F_{\mathrm{seg},\eta} - F_{\mathrm{diff},\eta} \,\Big] = 0\;}
    \label{eq:conservativa}
\end{equation}
con los tres flujos verticales
\begin{align}
    F_{\mathrm{adv},\eta} &= \Big(w - \eta\,\pd{h}{t} - \eta\,u\,\pd{h}{x}\Big)\phi_s \;=\; (h\,w_\eta)\,\phi_s , \label{eq:Fadv}\\
    F_{\mathrm{seg},\eta} &= -\,f_{sl}\,\phi_s(1-\phi_s), \label{eq:Fseg}\\
    F_{\mathrm{diff},\eta} &= \frac{D_{sl}}{h}\,\pd{\phi_s}{\eta}. \label{eq:Fdiff}
\end{align}
El término $w_\eta = \dfrac{1}{h}\big(w - \eta\,\partial_t h - \eta\,u\,\partial_x h\big)$
es la velocidad ``de malla'' que cruza las superficies $\eta=$ cte.

\subsection{Derivación detallada: de la Ec.~\eqref{eq:cart} a la Ec.~\eqref{eq:conservativa}}
\label{sec:derivacion}
Esta subsección desarrolla explícitamente el paso algebraico que conecta
la ley de conservación cartesiana \eqref{eq:cart} con la forma
conservativa en $\sigma$ \eqref{eq:conservativa}, en la notación de este
documento ($\eta = z/h(x,t)$, sin el desplazamiento $H_{\mathrm{base}}$
que usa la derivación alternativa de \texttt{dif\_adv\_seg.tex}).

\subsubsection{Paso 1: operadores diferenciales en las nuevas coordenadas}
Con $\eta = z/h(x,t)$, es decir $z = \eta\,h(x,t)$, la regla de la cadena
(manteniendo $z$ fijo al derivar respecto de $x$ o $t$) da:
\begin{align}
    \pd{\eta}{z} &= \frac{1}{h}
    &&\Longrightarrow\quad
    \boxed{\pd{}{z} = \frac{1}{h}\pd{}{\eta}}, \label{eq:dz}\\[6pt]
    \pd{\eta}{x}\bigg|_z &= -\frac{z}{h^2}\pd{h}{x} = -\frac{\eta}{h}\pd{h}{x}
    &&\Longrightarrow\quad
    \boxed{\pd{}{x}\bigg|_z = \pd{}{x}\bigg|_\eta - \frac{\eta}{h}\pd{h}{x}\,\pd{}{\eta}}, \label{eq:dx}\\[6pt]
    \pd{\eta}{t}\bigg|_z &= -\frac{\eta}{h}\pd{h}{t}
    &&\Longrightarrow\quad
    \boxed{\pd{}{t}\bigg|_z = \pd{}{t}\bigg|_\eta - \frac{\eta}{h}\pd{h}{t}\,\pd{}{\eta}}. \label{eq:dt}
\end{align}

\subsubsection{Paso 2: sustitución término a término en la Ec.~\eqref{eq:cart}}
Aplicando \eqref{eq:dz}--\eqref{eq:dt} a cada término de \eqref{eq:cart}:
\begin{align}
    \pd{\phi_s}{t} &= \pd{\phi_s}{t} - \frac{\eta}{h}\pd{h}{t}\pd{\phi_s}{\eta}, \\
    \pd{(u\phi_s)}{x} &= \pd{(u\phi_s)}{x} - \frac{\eta}{h}\pd{h}{x}\pd{(u\phi_s)}{\eta}, \\
    \pd{(w\phi_s)}{z} &= \frac1h\pd{(w\phi_s)}{\eta}, \qquad
    \pd{F_{\mathrm{seg}}}{z} = \frac1h\pd{F_{\mathrm{seg}}}{\eta}, \\
    \pd{}{z}\!\left(D_{sl}\pd{\phi_s}{z}\right) &= \frac1h\pd{}{\eta}\!\left(\frac{D_{sl}}{h}\pd{\phi_s}{\eta}\right).
\end{align}
Sumando todos los términos e igualando a cero, y multiplicando la
ecuación completa por $h$ para eliminar los denominadores, se obtiene:
\begin{multline}
    h\pd{\phi_s}{t} - \eta\,\pd{h}{t}\pd{\phi_s}{\eta}
    + h\pd{(u\phi_s)}{x} - \eta\,\pd{h}{x}\pd{(u\phi_s)}{\eta} \\
    + \pd{(w\phi_s)}{\eta} + \pd{F_{\mathrm{seg}}}{\eta}
    - \pd{}{\eta}\!\left(\frac{D_{sl}}{h}\pd{\phi_s}{\eta}\right) = 0.
    \label{eq:ast}
\end{multline}

\subsubsection{Paso 3: introducción de $Q=h\phi_s$ vía regla del producto}
Como $h$ no depende de $\eta$, la regla del producto da las identidades
\begin{align}
    \pd{Q}{t} = \pd{(h\phi_s)}{t} = h\pd{\phi_s}{t} + \phi_s\pd{h}{t}
    &\;\Longrightarrow\; h\pd{\phi_s}{t} = \pd{Q}{t} - \phi_s\pd{h}{t}, \label{eq:idQt}\\
    \pd{(uQ)}{x} = \pd{(hu\phi_s)}{x} = h\pd{(u\phi_s)}{x} + u\phi_s\pd{h}{x}
    &\;\Longrightarrow\; h\pd{(u\phi_s)}{x} = \pd{(uQ)}{x} - u\phi_s\pd{h}{x}. \label{eq:idQx}
\end{align}
Sustituyendo \eqref{eq:idQt} y \eqref{eq:idQx} en \eqref{eq:ast}:
\begin{multline}
    \left[\pd{Q}{t} - \phi_s\pd{h}{t}\right] - \eta\,\pd{h}{t}\pd{\phi_s}{\eta}
    + \left[\pd{(uQ)}{x} - u\phi_s\pd{h}{x}\right] - \eta\,\pd{h}{x}\pd{(u\phi_s)}{\eta} \\
    + \pd{(w\phi_s)}{\eta} + \pd{F_{\mathrm{seg}}}{\eta}
    - \pd{}{\eta}\!\left(\frac{D_{sl}}{h}\pd{\phi_s}{\eta}\right) = 0.
    \label{eq:paso3}
\end{multline}

\subsubsection{Paso 4: agrupación de los términos topográficos bajo $\partial/\partial\eta$}
Los términos que involucran $\partial h/\partial t$ en \eqref{eq:paso3} se
agrupan reconociendo la regla del producto inversa
$\partial(\eta\phi_s)/\partial\eta = \phi_s + \eta\,\partial\phi_s/\partial\eta$:
\begin{equation}
    -\phi_s\pd{h}{t} - \eta\,\pd{h}{t}\pd{\phi_s}{\eta}
    = -\pd{h}{t}\underbrace{\left[\phi_s + \eta\pd{\phi_s}{\eta}\right]}_{\partial(\eta\phi_s)/\partial\eta}
    = -\pd{}{\eta}\!\left(\eta\,\phi_s\,\pd{h}{t}\right),
    \label{eq:grupo_t}
\end{equation}
donde el último paso vale porque $\partial h/\partial t$ no depende de
$\eta$ y puede introducirse dentro de la derivada. Análogamente, con
$\partial(\eta u\phi_s)/\partial\eta = u\phi_s + \eta\,\partial(u\phi_s)/\partial\eta$:
\begin{equation}
    -u\phi_s\pd{h}{x} - \eta\,\pd{h}{x}\pd{(u\phi_s)}{\eta}
    = -\pd{}{\eta}\!\left(\eta\,u\,\phi_s\,\pd{h}{x}\right).
    \label{eq:grupo_x}
\end{equation}

\subsubsection{Paso 5: consolidación bajo un único operador $\partial/\partial\eta$}
Sustituyendo \eqref{eq:grupo_t} y \eqref{eq:grupo_x} en \eqref{eq:paso3},
todas las derivadas respecto a $\eta$ quedan bajo un solo corchete:
\begin{equation}
    \pd{Q}{t} + \pd{(uQ)}{x}
    + \pd{}{\eta}\left[\, w\phi_s - \eta\,\phi_s\pd{h}{t} - \eta\,u\,\phi_s\pd{h}{x}
    + F_{\mathrm{seg}} - \frac{D_{sl}}{h}\pd{\phi_s}{\eta} \,\right] = 0.
    \label{eq:paso5}
\end{equation}
Los tres primeros términos del corchete comparten el factor $\phi_s$ y se
factorizan directamente:
\begin{equation}
    w\phi_s - \eta\,\phi_s\pd{h}{t} - \eta\,u\,\phi_s\pd{h}{x}
    = \left(w - \eta\pd{h}{t} - \eta\,u\pd{h}{x}\right)\phi_s
    \equiv F_{\mathrm{adv},\eta},
\end{equation}
que es precisamente la definición \eqref{eq:Fadv}. Renombrando
$F_{\mathrm{seg}}\to F_{\mathrm{seg},\eta}$ y
$\dfrac{D_{sl}}{h}\dfrac{\partial\phi_s}{\partial\eta}\to F_{\mathrm{diff},\eta}$
en \eqref{eq:paso5}, se recupera exactamente la Ec.~\eqref{eq:conservativa}:
\begin{equation}
    \pd{Q}{t} + \pd{(uQ)}{x}
    + \pd{}{\eta}\Big[F_{\mathrm{adv},\eta} + F_{\mathrm{seg},\eta} - F_{\mathrm{diff},\eta}\Big] = 0. \qquad\blacksquare
\end{equation}

\paragraph{Nota sobre el signo de $F_{\mathrm{diff},\eta}$.} El signo
negativo delante de $F_{\mathrm{diff},\eta}$ en el corchete proviene de
que el término difusivo estaba originalmente del lado derecho de
\eqref{eq:cart} (restando, al pasarlo al lado izquierdo cambia de signo),
mientras que $F_{\mathrm{seg}}$ ya incluye su propio signo dentro de la
definición $F_{\mathrm{seg}}=-f_{sl}\,\phi_s(1-\phi_s)$.

% ======================================================================
\section{Campo de velocidades}
\label{sec:vel}
Para garantizar incompresibilidad exacta se derivan $u$ y $w$ de una
función de corriente $\psi = U_0 H_{\mathrm{base}}\,\eta^{m+1}\,g(x,t)$, con
$U_0$ la escala de advección, $m$ el exponente del perfil tipo Bagnold
(concentra el transporte cerca de $\eta=1$), y $g$ la modulación
espacio-temporal del \emph{grain-flow}:
\begin{equation}
    g(x,t) = 1 + A_{\mathrm{mod}}\,\sin\!\left(\frac{2\pi t}{T_{\mathrm{period}}}\right) S_x(x),
    \qquad
    S_x(x) = \tfrac12 + \tfrac12\tanh\!\left(\frac{x-x_{\mathrm{crest}}}{\delta_{\mathrm{cr}}}\right).
\end{equation}
\subsection{Velocidad horizontal (marco comóvil)}
\begin{equation}
    u(x,\eta,t) = \frac{U_0\,H_{\mathrm{base}}}{h(x)}\,(m+1)\,\eta^{m}\,g(x,t)\;-\;c_{\mathrm{mig}}.
    \label{eq:u}
\end{equation}
El término $-c_{\mathrm{mig}}$ es la recirculación del marco comóvil:
es el que produce el patrón de \emph{erosión en el stoss} y
\emph{deposición en el lee}, y hace que el interior enterrado avecte a
$\approx -c_{\mathrm{mig}}$.
\subsection{Velocidad vertical en $\sigma$ (libre de divergencia)}
Por continuidad, la velocidad en las caras $\eta$ (llamada
$\omega_{\mathrm{code}}$ en el código) es
\begin{equation}
    \omega_{\mathrm{code}}(x,\eta,t)
    = c_{\mathrm{mig}}\,\eta\,\frac{\dd h}{\dd x}
    - U_0\,H_{\mathrm{base}}\,\eta^{m+1}\,\pd{g}{x},
    \qquad
    \pd{g}{x} = A_{\mathrm{mod}}\sin\!\left(\tfrac{2\pi t}{T_{\mathrm{period}}}\right)\frac{\dd S_x}{\dd x}.
    \label{eq:wcode}
\end{equation}
Esta construcción satisface analíticamente
$\partial_x(h\,u) + \partial_\eta\,\omega_{\mathrm{code}} = 0$.

% ======================================================================
\section{Flujos físicos de Trewhela con capa activa confinada}
\label{sec:trewhela}
\label{sec:confin}
Los cierres para $f_{sl}$ y $D_{sl}$ siguen a Trewhela et al.\ (2021), pero
con dos modificaciones centrales de esta versión: el \textbf{desacople} de
la tasa de corte y su \textbf{confinamiento} a la capa activa.

\subsection{Tasa de corte desacoplada y confinada $\dot\gamma$}
En el modelo clásico, $\dot\gamma=|\partial u/\partial z|$ se calcula desde
la misma velocidad de advección \eqref{eq:u}, de modo que bajar $U_0$ para
suavizar la advección \emph{también} apaga la segregación. Aquí se
reconoce que el cizalle que impulsa el \emph{kinetic sieving} es el corte
\emph{interno} de la delgada capa móvil que avalancha, fijado por su
propia dinámica y no por la migración neta lenta del cuerpo. Se introduce
por ello una escala de corte independiente $U_{\mathrm{shear}}\gg U_0$:
\begin{equation}
    \dot\gamma(x,\eta,t) =
    \underbrace{\frac{U_{\mathrm{shear}}\,H_{\mathrm{base}}}{h_{\mathrm{eff}}^2(x)}\,m(m+1)\,\eta^{m-1}\,g(x,t)}_{\text{corte tipo Bagnold}}
    \;\cdot\;
    \underbrace{\vphantom{\frac{1}{h^2}}W_{\mathrm{act}}(x,\eta)}_{\text{confinamiento}},
    \label{eq:gammadot}
\end{equation}
donde $h_{\mathrm{eff}} = \max\!\big(h,\,2.5\,\delta_a\big)$ es un piso de
espesor que evita la divergencia $1/h^2$ en los pies delgados, y
$\delta_a=1.5$~mm es el espesor de la capa activa.

\paragraph{Ventana de capa activa.}
La segregación y la difusión granular solo ocurren en la capa superficial
que cizalla (espesor $\sim\delta_a$); debajo, los granos están bloqueados y
la estructura enterrada queda \emph{congelada}. Esto se impone con una
ventana tipo escalón suave en la profundidad física bajo la superficie,
$z_{\downarrow}=(1-\eta)\,h$:
\begin{equation}
    W_{\mathrm{act}}(x,\eta) = \tfrac12\left[\,1 - \tanh\!\left(\frac{z_{\downarrow}-\delta_a}{w_{\tanh}}\right)\right],
    \qquad w_{\tanh}=0.3~\text{mm}.
    \label{eq:Wact}
\end{equation}
Se eligió un escalón $\tanh$ y no una Gaussiana porque esta última todavía
vale $\sim6\%$ a media profundidad y, multiplicada por miles de sub-pasos
verticales, terminaba lavando las láminas enterradas. El escalón confina
$\dot\gamma$ (y por tanto $f_{sl}$ y $D_{sl}$) \emph{estrictamente} a la
capa activa.

\subsection{Diámetro medio, asimetría y velocidad de segregación}
\begin{align}
    \bar d(\phi_s) &= (1-\phi_s)\,d_l + \phi_s\,d_s, \\
    F(R,\phi_s) &= (R-1) + E\,(1-\phi_s)(R-1)^2, \\
    f_{sl} &= \frac{B\,\dot\gamma\,\bar d^2}{C\,\bar d + p(x,\eta)}\;F(R,\phi_s),
    \label{eq:fsl}\\
    D_{sl} &= A\,\dot\gamma\,\bar d^2 .
    \label{eq:Dsl}
\end{align}
El flujo segregativo completo es $F_{\mathrm{seg}}=-f_{sl}\,\phi_s(1-\phi_s)$
(el factor $\phi_s(1-\phi_s)$ apaga la segregación en zonas puras).

\paragraph{Piso de presión (regularización).}
La presión litostática granular es
\begin{equation}
    p(x,\eta) = \nu\,h(x)\,(1-\eta) \;+\; p_{\mathrm{floor}},
    \qquad p_{\mathrm{floor}} = 3\,\nu\,\delta_a ,
    \label{eq:p}
\end{equation}
con $\nu=0.6$ la fracción de empaquetamiento. En la superficie
($\eta\to1$) la presión litostática $\nu h(1-\eta)\to0$ y $f_{sl}$
\eqref{eq:fsl} divergía, forzando pasos de tiempo diminutos
($n_{\mathrm{sub}}\sim10^3$--$10^4$). El piso $p_{\mathrm{floor}}$
representa el \emph{overburden} mínimo bajo el cual el flujo de segregación
\emph{satura} en vez de diverger; reduce $f_{sl,\max}$ de $\sim$390 a
$\sim$18~mm/s y $n_{\mathrm{sub}}$ en un factor $\sim7$, disminuyendo así
la sobre-aplicación del operador vertical (y su difusión numérica).

\begin{table}[h!]
\centering
\caption{Constantes de Trewhela et al.\ (2021) y parámetros de la capa activa.}
\vspace{0.2cm}
\begin{tabular}{llcl}
\toprule
\textbf{Parámetro} & \textbf{Símbolo} & \textbf{Valor} & \textbf{Descripción}\\
\midrule
Constante de difusión    & $A$ & $0.108$ & Adimensional\\
Constante de segregación & $B$ & $0.3744$ & Adimensional\\
Constante de presión     & $C$ & $0.2712$ & Adimensional\\
Constante de asimetría   & $E$ & $2.0957$ & Adimensional\\
Empaquetamiento sólido   & $\nu$ & $0.6$ & Compactación del lecho\\
Escala de corte activo   & $U_{\mathrm{shear}}$ & $0.015$~m/s & Desacoplada de $U_0=3\times10^{-4}$~m/s\\
Espesor de capa activa   & $\delta_a$ & $1.5$~mm & Ancho de $W_{\mathrm{act}}$\\
Piso de presión          & $p_{\mathrm{floor}}$ & $3\nu\delta_a$ & Regulariza $f_{sl}$\\
\bottomrule
\end{tabular}
\end{table}

% ======================================================================
\section{Condiciones de borde: deposición por avalanchas}
\label{sec:borde}
\paragraph{Base ($\eta=0$).} Impermeable: $F_\eta|_{\eta=0}=0$.

\paragraph{Superficie ($\eta=1$).} Frontera abierta con balance de masa
según el signo de $\omega_{\mathrm{code}}$. En el \emph{stoss}
($\omega_{\mathrm{code}}\ge0$, erosión) el material sale con la
concentración local; en el \emph{lee} ($\omega_{\mathrm{code}}<0$,
deposición) entra con una composición sorteada
\begin{equation}
    \phi_{\mathrm{in}}(x,t) = \mathrm{clip}\!\Big(\lambda(t)\,P_{\mathrm{spat}}(x)\,P_{\mathrm{temp}}(t),\,0,\,1\Big),
\end{equation}
donde $P_{\mathrm{spat}}$ impone el sorteo espacial (grueso al pie, fino a
la cresta; Kleinhans 2004) y $\lambda(t)$ cierra el balance masa
erosionada $=$ masa depositada en cada instante:
\begin{equation}
    P_{\mathrm{spat}}(x) = \max\!\Big(0.1,\; 1 + 0.6\big(0.5 - s_{\mathrm{lee}}(x)\big)\Big),
    \qquad
    \lambda = \frac{\sum_{\mathrm{stoss}} \omega_{\mathrm{code}}\,\phi_s|_{\eta=1}}{\sum_{\mathrm{lee}} |\omega_{\mathrm{code}}|\,P_{\mathrm{spat}}\,P_{\mathrm{temp}}}.
\end{equation}

\paragraph{Avalanchas discretas (novedad).}
La modulación temporal que genera las láminas se agudiza respecto del seno
suave clásico para emular avalanchas \emph{discretas} de \emph{grain-flow}
(deposición concentrada en pulsos, con fuerte contraste grueso/fino):
\begin{equation}
    P_{\mathrm{temp}}(t) = 1 + A_P\,\operatorname{sign}(\sin\theta)\,|\sin\theta|^{\,p_{\mathrm{sharp}}},
    \qquad \theta = \frac{2\pi t}{T_{\mathrm{period}}},
    \label{eq:Ptemp}
\end{equation}
con $A_P=0.90$ (alta variabilidad) y $p_{\mathrm{sharp}}=0.5$ (pulso más
``cuadrado''). Cuanto mayor $A_P$ y menor $p_{\mathrm{sharp}}$, más nítida
la alternancia de láminas gruesas y finas.

\paragraph{Bordes horizontales.} Dominio periódico en $x$: el material que
sale por un extremo reingresa por el otro, cerrando el ciclo de reciclaje
(grueso al pie del lee $\to$ enterramiento $\to$ reexposición en el stoss).

% ======================================================================
\section{Condición inicial}
\label{sec:ci}
La duna arranca como una \textbf{mezcla homogénea},
\begin{equation}
    \phi_s(x,\eta,0) = \phi_{s,\mathrm{target}} = 0.70 \quad \text{en todo el dominio},
\end{equation}
sin gradación ni estructura sembrada. Todo el patrón interno (armadura de
gruesos superficial, lag basal de gruesos, láminas) \emph{emerge} de la
dinámica de segregación confinada más la deposición por avalanchas.

% ======================================================================
\section{Método numérico}
\label{sec:numerico}
Se usa \emph{operator splitting}: un paso horizontal grande $\Delta t_h$
seguido de un sub-ciclo vertical de $n_{\mathrm{sub}}$ pasos. Los
coeficientes que dependen de $t$ (pero no de $\phi_s$) se ``congelan''
durante el sub-ciclo, evaluados en $t+\tfrac12\Delta t_h$, lo que es válido
porque $\Delta t_h\ll T_{\mathrm{period}}$.

\subsection{Paso horizontal: SSP-RK2 con reconstrucción WENO5}
El transporte horizontal $-\partial_x(uQ)$ se integra con Runge--Kutta de
2.\textsuperscript{o} orden de estabilidad fuerte:
\begin{align}
    Q^\ast &= Q^n + \Delta t_h\,\mathcal{R}_x(Q^n,t^n),\\
    Q^{n+1} &= \tfrac12 Q^n + \tfrac12\Big(Q^\ast + \Delta t_h\,\mathcal{R}_x(Q^\ast,t^n{+}\Delta t_h)\Big).
\end{align}
Los estados en las caras se reconstruyen con \textbf{WENO5}
(5.\textsuperscript{o} orden). Para la cara $i+\tfrac12$ se combinan tres
\emph{stencils} de tres celdas mediante pesos no lineales $\omega_k$
basados en indicadores de suavidez $\beta_k$:
\begin{equation}
    \phi^{L}_{i+1/2} = \sum_{k=0}^{2}\omega_k\,p_k,
    \qquad
    \omega_k = \frac{\alpha_k}{\sum_l \alpha_l},
    \quad
    \alpha_k = \frac{d_k}{(\varepsilon+\beta_k)^2},
\end{equation}
con pesos lineales $d=(0.1,\,0.6,\,0.3)$ y $\varepsilon=10^{-6}$. En zonas
suaves WENO5 alcanza 5.\textsuperscript{o} orden (difusión numérica
$\sim\Delta x^5$), y cerca de un frente se reduce a un \emph{stencil} ENO
no oscilatorio. Este esquema reemplazó al MUSCL de 2.\textsuperscript{o}
orden ($\sim\Delta x^3$) precisamente porque este último difuminaba las
láminas de foreset advectadas. Como el flujo $u\phi$ es lineal en $\phi$
(para $u$ dado), basta el flujo \emph{upwind} exacto con los estados
reconstruidos.

\subsection{Sub-ciclo vertical: MUSCL + Rusanov + difusión centrada}
Cada sub-paso vertical aplica Euler explícito al flujo
$F_\eta=F_{\mathrm{adv},\eta}+F_{\mathrm{seg},\eta}-F_{\mathrm{diff},\eta}$.
Los estados en las caras interiores se reconstruyen con pendientes
limitadas por \textbf{minmod}:
\begin{equation}
    \sigma_j = \operatorname{minmod}\!\big(\phi_j-\phi_{j-1},\ \phi_{j+1}-\phi_j\big),
    \qquad
    \phi^{\pm} = \phi_j \pm \tfrac12\sigma_j .
\end{equation}
\begin{itemize}
    \item \textbf{Advección vertical} (lineal): \emph{upwind} exacto con los
    estados reconstruidos.
    \item \textbf{Segregación} (no lineal, $F(\phi)=-f_{sl}\phi(1-\phi)$):
    flujo de \textbf{Rusanov} (Lax--Friedrichs local)
    \begin{equation}
        F^{\mathrm{num}}_{\mathrm{seg}} = \tfrac12\big[F(\phi^L)+F(\phi^R)\big]
        - \tfrac12\,a_{\max}\,(\phi^R-\phi^L),
        \quad
        a_{\max} = f_{sl}\,\max\!\big(|1-2\phi^L|,\,|1-2\phi^R|\big),
    \end{equation}
    donde $a_{\max}$ acota $|\dd F/\dd\phi|$.
    \item \textbf{Difusión}: diferencias centradas (ya es 2.\textsuperscript{o}
    orden) del flujo $F_{\mathrm{diff},\eta}=(D_{sl}/h)\,\partial_\eta\phi$.
\end{itemize}

\subsection{Paso de tiempo (CFL) y sub-ciclado}
\begin{equation}
    \Delta t_h = 0.5\,\frac{\dd x}{u_{\max}},
    \qquad
    \Delta t_v = \min\!\left(
    0.5\,\frac{\dd\eta\,\min h}{w_{\eta,\max}+f_{sl,\max}},\;\;
    0.45\,\frac{(\dd\eta\,\min h)^2}{D_{sl,\max}}\right),
\end{equation}
y $n_{\mathrm{sub}}=\lceil \Delta t_h/\Delta t_v\rceil$. El primer término
de $\Delta t_v$ es la restricción advectiva/segregativa y el segundo la
difusiva (parabólica).

\subsection{Acotamiento y conservación de masa}
Tras cada operación se impone $\phi_s\in[0,1]$: un recorte simple dentro
del sub-ciclo y un recorte \emph{conservativo} (\texttt{\_clip\_conserv},
barrido bidireccional en $\eta$ que redistribuye el exceso/déficit dentro
de cada columna) al final de cada paso. Finalmente se aplica una
corrección multiplicativa global
\begin{equation}
    Q \leftarrow Q\cdot \frac{M_0}{\sum_{i,j}Q_{i,j}\,\dd x\,\dd\eta},
\end{equation}
que elimina el error de redondeo acumulado. Con esto,
$\langle\phi_s\rangle=0.700000$ se mantiene constante en toda la corrida.

\paragraph{Malla y tiempo.} $N_x=160$, $N_z=40$, $t_{\max}=2400$~s. La
alta resolución, junto con WENO5 y el piso de presión, es lo que permite
que las láminas advectadas sobrevivan sin lavarse numéricamente.

% ======================================================================
\section{Adaptaciones para la versión curva}
\label{sec:curva}
La versión curva (\texttt{dune\_bidispersa\_deposicion\_ext\_curva.py})
comparte \emph{exactamente} toda la física y numérica anteriores
(Secciones \ref{sec:cambio}--\ref{sec:numerico}). El \textbf{único} cambio
es la geometría $h(x)$, para pasar de la duna triangular a una silueta de
duna real (comparable con la imagen experimental). Se introdujeron tres
modificaciones sobre \eqref{eq:h_tri}:

\paragraph{1. Stoss convexo (ley de potencia).} El barlovento recto se
reemplaza por una ley de potencia con exponente $p_{\mathrm{stoss}}<1$, que
lo hace \emph{convexo} (se empina rápido desde el pie y se aplana hacia la
cresta, abombándose por encima de la recta, $\dd^2h/\dd x^2<0$):
\begin{equation}
    h_{\mathrm{stoss}}(x) = H_{\mathrm{base}} + H_d\left(\frac{x}{x_{\mathrm{crest}}}\right)^{p_{\mathrm{stoss}}},
    \qquad p_{\mathrm{stoss}} = 0.55 .
\end{equation}
El lee se mantiene recto al ángulo de reposo (30\textsuperscript{o}).

\paragraph{2. Suavizado de la cresta (continuidad $C^1$).} El perfil a
trozos $h_{\mathrm{tent}}$ (stoss convexo $+$ lee recto) se convoluciona
con un filtro Gaussiano periódico para eliminar el quiebre de pendiente en
la cresta:
\begin{equation}
    \tilde h = \mathcal{G}_{\sigma_s} \ast h_{\mathrm{tent}},
    \qquad \sigma_s = 6~\text{celdas} \ (\approx 3.75~\text{mm}).
\end{equation}
Como el núcleo Gaussiano es no negativo, el suavizado es un promedio
ponderado de pendientes vecinas y \emph{nunca} excede la pendiente máxima
original; así $\alpha_{\mathrm{lee}}=30^\circ$ sigue siendo la pendiente
\emph{máxima} local del lee (ya no constante).

\paragraph{3. Reescalado y pendiente continua.} Se reescala para conservar
la altura física $H_d$ en la cresta y se calcula $\dd h/\dd x$ por
diferencias centradas periódicas:
\begin{equation}
    h = H_{\mathrm{base}} + (\tilde h - H_{\mathrm{base}})\,\frac{H_d}{\max(\tilde h)-H_{\mathrm{base}}},
    \qquad
    \frac{\dd h}{\dd x}\Big|_i = \frac{h_{i+1}-h_{i-1}}{2\,\dd x}.
\end{equation}

\paragraph{Efecto físico.} La motivación numérica es que la topografía
triangular tiene $\dd h/\dd x$ \emph{discontinua} en la cresta, y esa
discontinuidad entra directamente en la velocidad vertical
$\omega_{\mathrm{code}}$ \eqref{eq:wcode} (vía $c_{\mathrm{mig}}\,\eta\,
\dd h/\dd x$), contaminando el corte de segregación cerca de la cresta. Con
la geometría curva $\dd h/\dd x$ es continua en toda la duna: desaparece
ese artefacto y la silueta (cresta redondeada $+$ stoss convexo abombado)
reproduce la forma observada en el experimento. La estructura de
segregación resultante (lag basal de gruesos bajo un cuerpo rico en finos)
es la misma que en la versión triangular, confirmando que la forma altera
la apariencia pero no el mecanismo de \emph{sorting}.

\end{document}
